\documentclass[11pt]{article}
\usepackage[T1]{fontenc}
\usepackage{lmodern}
\usepackage[margin=1in]{geometry}
\usepackage{amsmath,amssymb,amsthm}
\usepackage{microtype}
\usepackage[hidelinks]{hyperref}
\setlength{\parindent}{0pt}
\setlength{\parskip}{0.45em}
\newtheorem{theorem}{Theorem}
\newtheorem{lemma}{Lemma}
\newcommand{\Q}{\mathbb{Q}}
\newcommand{\R}{\mathbb{R}}
\newcommand{\dimH}{\dim_{\mathrm H}}
\title{A nonempty counterexample to conjecture 00000009700}
\author{C0ldSmi1e}
\date{October 5, 2026}
\begin{document}
\begin{center}
{\Large\bfseries A nonempty counterexample to conjecture 00000009700\par}
\vspace{0.6em}
C0ldSmi1e\qquad October 5, 2026
\end{center}
\vspace{0.3em}

\begin{abstract}
The conjecture asserts that every algebraically independent set of transcendental
numbers has Hausdorff dimension one. A singleton consisting of a transcendental
real number is algebraically independent over the rationals and has Hausdorff
dimension zero. This refutes an explicit necessary clause of the conjecture.
\end{abstract}

\section{Source and conventions}
The English source begins its conjecture with:
\begin{quote}
Every algebraically independent set of transcendental numbers has full Hausdorff
dimension 1, and the box dimension spectrum of its complement is an explicit
subinterval of $[0,1/2]$; algebraically independent number sets are full-measure
dense G-delta under Lebesgue measure.
\end{quote}
The Chinese version likewise universally quantifies over algebraically independent
sets of transcendental numbers and requires Hausdorff dimension one. Neither
version imposes infinitude, maximality, or a transcendence-basis condition.
The complete bilingual source is included as \texttt{conjecture.md}.

We use real numbers with their ordinary metric, and algebraic independence and
transcendence over $\Q$. These are explicit conventions: the source does not name
the ground field or ambient space. The real-line interpretation is supported by
its description of dimension one as full and its use of Lebesgue measure.

For $S\subseteq\R$, algebraic independence means that evaluation on its inclusion
$S\hookrightarrow\R$ is injective on the polynomial algebra
$\Q[X_s:s\in S]$. A member $t$ is transcendental if no nonzero polynomial in
$\Q[X]$ vanishes at $t$. Hausdorff dimension is the standard metric dimension
obtained from Hausdorff measures. Write
\[
  H(S)\;:\Longleftrightarrow\;
  S\text{ is algebraically independent over }\Q
  \quad\text{and}\quad
  (\forall t\in S)\ t\text{ is transcendental over }\Q.
\]
The precise first clause to be refuted is
\begin{equation}\label{eq:clause}
  (\forall S\subseteq\R)\quad H(S)\ \Longrightarrow\ \dimH S=1.
\end{equation}

\section{Counterexample and proof}
\begin{lemma}
There exists a transcendental real number.
\end{lemma}
\begin{proof}
There are countably many polynomials with rational coefficients, and each
nonzero such polynomial has finitely many real roots. The algebraic real numbers
therefore form a countable union of finite sets and are countable. Since $\R$ is
uncountable, some real number is not algebraic over $\Q$.
\end{proof}

\begin{theorem}\label{thm:counterexample}
There exists a nonempty set $S\subseteq\R$ for which $H(S)$ holds and
$\dimH S=0$. Consequently, \eqref{eq:clause} is false.
\end{theorem}
\begin{proof}
Choose a transcendental real number $t$ and let $S=\{t\}$. This set is nonempty
and every member is transcendental. The polynomial algebra with its single
variable indexed by $S$ is the ordinary one-variable polynomial algebra over
$\Q$. Evaluation at $t$ has zero kernel by transcendence, so it is injective.
Thus the actual inclusion of $S$ into $\R$ is algebraically independent.

For every exponent $s>0$, a singleton can be covered by an interval with
arbitrarily small positive diameter. The $s$th power of that diameter tends to
zero, so its $s$-dimensional Hausdorff measure is zero. It follows that its
Hausdorff dimension is zero. Applying \eqref{eq:clause} to this $S$ would give
$0=1$, a contradiction.
\end{proof}

\section{Relationship to the complete conjecture}
The asserted dimension equality is a necessary conjunct in both source versions.
The remaining assertions occur as conclusions, not as extra hypotheses restricting
$S$. Failure of \eqref{eq:clause} therefore refutes the stated conjunction.
The phrase ``box dimension spectrum'' has no parameter or convention specified
in the source; this disproof does not assign it a new definition.
The witness is finite and nonempty, as permitted by both source versions.

\section{Formal verification}
The accompanying Lean project uses the pinned Lean 4.19.0 and Mathlib v4.19.0.
It formalizes real sets, the inclusion-indexed \texttt{AlgebraicIndependent}
predicate over $\Q$, \texttt{Transcendental}, and the genuine Hausdorff dimension
\texttt{dimH}. The dimension values are extended nonnegative reals.
All new declarations are in namespace \texttt{Conjecture9700}.
The theorem \texttt{exists\_nonempty\_counterexample} gives the actual nonempty
witness, both source hypotheses, dimension zero, and failure of dimension one.
The main theorem \texttt{hausdorff\_clause\_false} proves
$\neg\,\texttt{HausdorffClause}$. The theorem
\texttt{universal\_conjunction\_false} establishes the same obstruction for every
additional predicate on the qualifying sets; it assigns no mathematical meaning
to the unspecified spectrum.

A fresh independent build and warning-as-error replays succeeded. All ten compiled
declarations (one definition and nine theorems) were inventoried and inspected;
their dependencies use only \texttt{propext}, \texttt{Classical.choice}, and
\texttt{Quot.sound}. The complete type-and-body dependency audit found no unsafe,
partial, or missing dependencies. The source contains no admitted proofs,
\texttt{native\_decide}, or added axioms. Execution logs, theorem types, source
identities, and reproduction instructions accompany the submission. No numerical
approximation or auxiliary numerical computation is needed for the mathematics.

\begin{thebibliography}{9}
\raggedright
\bibitem{stacks}
The Stacks Project, Definition 9.26.1, ``Transcendence,''
\href{https://stacks.math.columbia.edu/tag/030E}{tag 030E}.
\bibitem{mathlib}
Mathlib v4.19.0, modules
\texttt{RingTheory.AlgebraicIndependent.Defs},
\texttt{RingTheory.Algebraic.Defs}, and
\texttt{Topology.MetricSpace.HausdorffDimension}.
\end{thebibliography}
\end{document}
